% Source printed pp. 165-175 / PDF pp. 173-183.
\section{拟凝聚层与上同调}\label{src:5.1}
\begin{epigraph}
若不受牢固的人际关系制约, 以抽象观念为指导是件危险的事.
\epigraphauthor{Paul Feyerabend, \textit{科学的暴政}}
\end{epigraph}

\subsection{层}\label{src:5.1.1-subsection}
Deligne--Mumford 叠的小平展景可类似于概形的小平展景定义 (例~\bookref{3.2.4}).
\begin{definition}\label{src:5.1.1}
Deligne--Mumford 叠 $\mathcal X$ 的\newterm{小平展景}{small étale site} 是由 $\mathcal X$ 上的平展概形组成的范畴 $\mathcal X_{\text{ét}}$. 一个 $\mathcal X$-概形 $U$ 的覆盖是 $\mathcal X$ 上的一族平展态射 $\{U_i\to U\}$, 使得 $\coprod_i U_i\to U$ 满射.
\end{definition}
于是可以讨论 $\mathcal X_{\text{ét}}$ 上的 Abel 群层及其态射. 以 $\mathsf{Ab}(\mathcal X_{\text{ét}})$ 记 $\mathcal X_{\text{ét}}$ 上 Abel 层的范畴. 对 $\mathcal X_{\text{ét}}$ 上的 Abel 层 $F$, 以 $F(U\to\mathcal X)$ 记它在平展 $\mathcal X$-概形 $U$ 上的截面.
\begin{example}[结构层]\label{src:5.1.2}
Deligne--Mumford 叠上的\newterm{结构层}{structure sheaf} $\mathcal O_{\mathcal X}$ 在平展 $\mathcal X$-概形 $U$ 上定义为 $\mathcal O_{\mathcal X}(U\to\mathcal X)=\Gamma(U,\mathcal O_U)$.
\end{example}
\begin{example}[微分]\label{src:5.1.3}
若 $\mathcal X$ 是概形 $S$ 上的 Deligne--Mumford 叠, 则\newterm{相对微分层}{relative sheaf of differentials} $\Omega_{\mathcal X/S}$ 定义为 $\Omega_{\mathcal X/S}(U\to\mathcal X)=\Gamma(U,\Omega_{U/S})$.
\end{example}
\begin{example}[Hodge 丛]\label{src:5.1.4}
在 $\mathcal M_g$ ($g\geq2$) 上如下定义层 $\mathcal H$: 对由光滑曲线族 $\mathcal C\to U$ 分类的平展态射 $U\to\mathcal M_g$, 置 $\mathcal H(U\to\mathcal M_g)=\Gamma(\mathcal C,\Omega_{\mathcal C/U})$. 稍后将看到 $\mathcal H$ 是凝聚 $\mathcal O_{\mathcal M_g}$-模, 事实上是秩 $g$ 的\newterm{Hodge 丛}{Hodge bundle}.
\end{example}
\begin{exercise}[简单]\label{src:5.1.5}
若 $k$ 是代数闭域, 证明由 $F\mapsto F(\Spec k)$ 定义的函子 $\mathsf{Ab}((\Spec k)_{\text{ét}})\to\mathsf{Ab}$ 是范畴等价.
\end{exercise}
虽然 $\mathcal X_{\text{ét}}$ 上的层 $F$ 的截面只在平展 $\mathcal X$-概形上定义, 但可如下将定义延伸到 $\mathcal X$ 上的平展 Deligne--Mumford 叠 $\mathcal U$: 选取来自概形的平展呈示 $U\to\mathcal U$ 和 $R\to U\times_{\mathcal U}U$, 并定义
\[
F(\mathcal U\to\mathcal X):=\operatorname{Eq}\bigl(F(U\to\mathcal X)\rightrightarrows F(R\to\mathcal X)\bigr).
\]
可以检验这与呈示的选择无关. 特别地, 可以讨论整体截面
\[
\Gamma(\mathcal X,F):=F(\mathcal X\xrightarrow{\mathrm{id}}\mathcal X).
\]
\begin{exercise}[简单]\label{src:5.1.6}
若 $F$ 是 Deligne--Mumford 叠 $\mathcal X$ 上的 Abel 层, 证明 $\Gamma(\mathcal X,F)=\operatorname{Hom}_{\mathsf{Ab}(\mathcal X_{\text{ét}})}(\underline{\mathbb Z},F)$, 其中 $\underline{\mathbb Z}$ 是常值层.
\end{exercise}
给定 Deligne--Mumford 叠的态射 $f:\mathcal X\to\mathcal Y$, 有函子
\[
\begin{tikzcd}
\mathsf{Ab}(\mathcal X_{\text{ét}})\arrow[r,shift left,"f_*"]&\mathsf{Ab}(\mathcal Y_{\text{ét}})\arrow[l,shift left,"f^{-1}"] .
\end{tikzcd}
\]
其中\newterm{推前}{pushforward} 定义为 $f_*F(V\to\mathcal Y):=F(\mathcal X\times_{\mathcal Y}V\to\mathcal X)$, 而\newterm{逆像}{inverse image} $f^{-1}G$ 是预层
\[
(U\to\mathcal X)\longmapsto\mathop{\mathrm{colim}}_{(V\to\mathcal Y,\,U\to\mathcal X\times_{\mathcal Y}V)}G(V\to\mathcal Y)
\]
的层化. 此处余极限遍历平展态射对 $V\to\mathcal Y$ 与 $U\to\mathcal X\times_{\mathcal Y}V$, 即交换图
\[
\begin{tikzcd}
U\arrow[r]\arrow[d,"\text{平展}"']&V\arrow[d,"\text{平展}"]\\
\mathcal X\arrow[r]&\mathcal Y.
\end{tikzcd}
\]
\begin{exercise}[简单]\label{src:5.1.7}
证明若 $f:\mathcal X\to\mathcal Y$ 是 Deligne--Mumford 叠的平展态射, $G$ 是 $\mathcal Y_{\text{ét}}$ 上的层, 则 $f^{-1}G=G|_{\mathcal X}$ 是 $G$ 在 $\mathcal X_{\text{ét}}$ 上的限制, 即层 $(U\to\mathcal X)\mapsto G(U\to\mathcal X\to\mathcal Y)$. 特别地, 证明 $f^{-1}\mathcal O_{\mathcal Y}=\mathcal O_{\mathcal X}$.
\end{exercise}
\begin{exercise}[简单]\label{src:5.1.8}
设 $X$ 是概形, $x\in X$.
\begin{enumerate}[(a)]
\item 记 $i_x:\Spec\kappa(x)\to X$, 证明 $i_x^{-1}\mathcal O_X$ 的整体截面是 $\mathcal O_{X,x}$ 的 Hensel 化 $\mathcal O^h_{X,x}$.
\item 记 $i_{\overline x}:\Spec\overline{\kappa(x)}\to X$, 证明 $i_{\overline x}^{-1}$ 可与 $\mathcal O_{X,x}$ 的严格 Hensel 化 $\mathcal O^{\mathrm{sh}}_{X,x}$ 等同.
\end{enumerate}
关于 (严格) Hensel 化的背景, 见 \S\bookref{B.5.3-subsection}.
\end{exercise}
\begin{exercise}\label{src:5.1.9}
设 $k$ 是域, $k^{\mathrm{sep}}$ 是可分闭包, $G:=\operatorname{Gal}(k^{\mathrm{sep}}/k)$ 是赋予 profinite 拓扑的 Galois 群.
\begin{enumerate}[(a)]
\item 证明协变函子 $F:\{\text{平展 $k$-代数}\}\to\mathsf{Ab}$ 定义 $(\Spec k)_{\text{ét}}$ 上的 Abel 群层, 当且仅当对每个有限平展 $k$-代数族 $\{A_i\}$ 有 $F(\prod_i A_i)=\bigoplus_i F(A_i)$, 且对 $k$ 的有限可分域扩张之间的每个有限 Galois 扩张 $L/K$ 有 $F(K)=F(L)^{\operatorname{Gal}(L/K)}$.
\item 证明 $(\Spec k)_{\text{ét}}$ 上层的范畴等价于离散 $G$-模的范畴, 即带连续 $G$-作用的 Abel 群的范畴.

\textit{提示:} 给定 $F$, 定义离散 $G$-模 $M_F:=\mathrm{colim}_{k'/k}F(k')$, 其中余极限遍历有限可分域扩张. 反之, 给定 $M$, 证明有限域扩张上的函子 $k'\mapsto M^{\operatorname{Gal}(k^{\mathrm{sep}}/k')}$ 延伸为 $(\Spec k)_{\text{ét}}$ 上的层 $F_M$. 证明 $M\mapsto M_F$ 与 $M\mapsto F_M$ 是互逆函子.
\end{enumerate}
\end{exercise}
\newpage
\begin{exercise}\label{src:5.1.10}
设 $\mathcal X$ 是 Deligne--Mumford 叠, $x:\Spec k\to\mathcal X$ 是几何点 (即 $k$ 是代数闭域). 定义 $\mathcal X_{\text{ét}}$ 上的 Abel 层 $F$ 在 $x$ 处的\newterm{茎}{stalk} 为 $F_x:=x^{-1}F$.
\begin{enumerate}[(a)]
\item 证明由 $F\mapsto F_x$ 定义的函子 $\mathsf{Ab}(\mathcal X_{\text{ét}})\to\mathsf{Ab}$ 正合.
\item 证明 $\mathcal X_{\text{ét}}$ 上的 Abel 层序列正合, 当且仅当对每个几何点, 所诱导的茎序列正合.
\item 证明对 Deligne--Mumford 叠的态射 $f:\mathcal X\to\mathcal Y$, $f^{-1}$ 正合.
\end{enumerate}
\end{exercise}
\begin{exercise}[细节]\label{src:5.1.11}
证明 $f^{-1}$ 是 $f_*$ 的左伴随.
\end{exercise}
\begin{exercise}[技术性]\label{src:5.1.12}
若 $\mathcal X$ 是 Deligne--Mumford 叠, 另定义景 $\mathcal X_{\text{ét}'}$ 为 $\mathcal X$ 上平展代数空间的范畴, 覆盖取平展覆盖. 证明 $\mathcal X_{\text{ét}}$ 与 $\mathcal X_{\text{ét}'}$ 上层的范畴等价.
\end{exercise}

\subsection{\texorpdfstring{$\mathcal O_{\mathcal X}$}{OX}-模}\label{src:5.1.2-subsection}
结构层 $\mathcal O_{\mathcal X}$ 是 $\mathsf{Ab}(\mathcal X_{\text{ét}})$ 中的环对象, $\mathcal O_{\mathcal X}$-模定义为 $\mathcal O_{\mathcal X}$ 上的模对象.
\begin{definition}\label{src:5.1.13}
若 $\mathcal X$ 是 Deligne--Mumford 叠, 则\newterm{$\mathcal O_{\mathcal X}$-模层}{sheaf of $\mathcal O_{\mathcal X}$-modules} (或简称 $\mathcal O_{\mathcal X}$-模) 是 $\mathcal X_{\text{ét}}$ 上的层 $F$, 它是层范畴中 $\mathcal O_{\mathcal X}$ 的模对象, 即对每个平展 $\mathcal X$-概形 $U$, $F(U\to\mathcal X)$ 是 $\mathcal O_{\mathcal X}(U\to\mathcal X)=\Gamma(U,\mathcal O_U)$ 上的模, 且模结构与沿平展 $\mathcal X$-概形之间态射 $V\to U$ 的限制相容.
\end{definition}
以 $\Mod(\mathcal O_{\mathcal X})$ 记 $\mathcal O_{\mathcal X}$-模的范畴. 给定两个 $\mathcal O_{\mathcal X}$-模 $F,G$, 可将\newterm{张量积}{tensor product} $F\otimes G:=F\otimes_{\mathcal O_{\mathcal X}}G$ 定义为由 $(U\to\mathcal X)\mapsto F(U\to\mathcal X)\otimes_{\mathcal O_{\mathcal X}(U\to\mathcal X)}G(U\to\mathcal X)$ 给出的 $\mathcal O_{\mathcal X}$-模的层化. 对来自概形的平展态射 $f:U\to\mathcal X$, \newterm{Hom 层}{Hom sheaf} $\mathcal H\!om_{\mathcal O_{\mathcal X}}(F,G)$ 在其上的截面为 $\operatorname{Hom}_{\mathcal O_U}(F|_U,G|_U)$, 其中 $F|_U=f^{-1}F$ 表示 $F$ 在 $U_{\text{ét}}$ 上的限制.
\begin{exercise}[简单]\label{src:5.1.14}
若 $F$ 是 $\mathcal O_{\mathcal X}$-模, 证明 $\Gamma(\mathcal X,F)=\operatorname{Hom}_{\mathcal O_{\mathcal X}}(\mathcal O_{\mathcal X},F)$.
\end{exercise}
给定 Deligne--Mumford 叠的态射 $f:\mathcal X\to\mathcal Y$, 有函子
\[
\begin{tikzcd}
\Mod(\mathcal O_{\mathcal X})\arrow[r,shift left,"f_*"]&\Mod(\mathcal O_{\mathcal Y})\arrow[l,shift left,"f^*"] .
\end{tikzcd}
\]
其中对 $\mathcal O_{\mathcal X}$-模 $F$, $f_*F$ 是作为层的推前, 并自然成为 $\mathcal O_{\mathcal Y}$-模. 对 $\mathcal O_{\mathcal Y}$-模 $G$, 由于在 $\mathcal X_{\text{ét}}$ 上有环层态射 $f^{-1}\mathcal O_{\mathcal Y}\to\mathcal O_{\mathcal X}$, 且 $f^{-1}G$ 是 $f^{-1}\mathcal O_{\mathcal Y}$-模, 可以将\newterm{拉回}{pullback} 定义为 $\mathcal O_{\mathcal X}$-模
\[
f^*G:=f^{-1}G\otimes_{f^{-1}\mathcal O_{\mathcal Y}}\mathcal O_{\mathcal X}.
\]
若 $f$ 平展, 则 $f^*G=f^{-1}G=G|_{\mathcal X}$, 因为 $f^{-1}\mathcal O_{\mathcal Y}=\mathcal O_{\mathcal X}$ (习题~\bookref{5.1.7}).
\begin{exercise}[细节]\label{src:5.1.15}
证明 $f^*$ 是 $f_*$ 的左伴随.
\end{exercise}
\begin{exercise}[细节]\label{src:5.1.16}
证明 $\Mod(\mathcal O_{\mathcal X})$ 是 Abel 范畴.
\end{exercise}

\subsection{拟凝聚层}\label{src:5.1.3-subsection}
设 $F$ 是 Deligne--Mumford 叠 $\mathcal X$ 上的 $\mathcal O_{\mathcal X}$-模. 对平展 $\mathcal X$-概形 $U$, 有 $F$ 在 $U$ 的小平展景 $U_{\text{ét}}$ 上的限制 $F|_{U_{\text{ét}}}$, 以及进一步在 $U$ 的 Zariski 拓扑上的限制 $F|_{U_{\mathrm{Zar}}}$. 注意当 $X$ 是概形时, $\mathcal O_X$ 既可指 $X_{\text{ét}}$ 上的结构层, 也可指 $X_{\mathrm{Zar}}$ 上的结构层; 可能混淆时, 分别写作 $\mathcal O_{X_{\text{ét}}}$ 和 $\mathcal O_{X_{\mathrm{Zar}}}$.
\begin{definition}\label{src:5.1.17}
设 $\mathcal X$ 是 Deligne--Mumford 叠. 一个 $\mathcal O_{\mathcal X}$-模 $F$ 称为\newterm{拟凝聚的}{quasi-coherent}, 若
\begin{enumerate}[(1)]
\item 对每个平展 $\mathcal X$-概形 $U$, 限制 $F|_{U_{\mathrm{Zar}}}$ 是拟凝聚 $\mathcal O_{U_{\mathrm{Zar}}}$-模, 且
\item 对平展 $\mathcal X$-概形之间的每个平展态射 $f:U\to V$, 自然态射 $f^*(F|_{V_{\mathrm{Zar}}})\to F|_{U_{\mathrm{Zar}}}$ 是同构.
\end{enumerate}
$\mathcal X$ 上的拟凝聚 $\mathcal O_{\mathcal X}$-模 $F$ 称为\newterm{向量丛}{vector bundle} (分别为秩 $r$ 的向量丛, 线丛), 若对来自概形的每个平展态射 $U\to\mathcal X$, $F|_{U_{\mathrm{Zar}}}$ 都具有相应性质. 若此外 $\mathcal X$ 局部 Noether, 则称 $F$ 为\newterm{凝聚的}{coherent}, 若对每个平展态射 $U\to\mathcal X$, $F|_{U_{\mathrm{Zar}}}$ 凝聚.
\end{definition}
以 $\QCoh(\mathcal X)$ 和 $\operatorname{Coh}(\mathcal X)$ (在 Noether 情形) 分别记拟凝聚层和凝聚层的范畴. 稍后将看到当 $\mathcal X$ 是概形时, 这与通常的定义一致; 先给出几个例子.
\begin{example}\label{src:5.1.18}
结构层 $\mathcal O_{\mathcal X}$ 拟凝聚. 它是线丛, 当 $\mathcal X$ 局部 Noether 时凝聚.
\end{example}
\begin{example}\label{src:5.1.19}
对概形 $S$ 上的 Deligne--Mumford 叠 $\mathcal X$, 例~\bookref{5.1.3} 的相对微分层 $\Omega_{\mathcal X/S}$ 拟凝聚, 因为对平展 $\mathcal X$-概形之间的平展态射 $f:U\to V$, $f^*\Omega_{V/S}\to\Omega_{U/S}$ 是同构. 当 $\mathcal X\to S$ 光滑时, 它是向量丛.
\end{example}
\begin{example}\label{src:5.1.20}
对 $\mathcal M_g$ ($g\geq2$), 例~\bookref{5.1.4} 的 Hodge 丛 $\mathcal H$ 是秩 $g$ 的向量丛. 这由命题~\bookref{6.1.16}(2) 推出: 对应于 $\mathcal M_g$-概形 $V$ 的亏格 $g$ 光滑曲线族 $\pi:\mathcal C\to V$, 构造 $\pi_*\Omega_{\mathcal C/V}$ 与沿映射 $f:U\to V$ 的基变换交换, 即
\[
f^*(\pi_*\Omega_{\mathcal C/V})\xrightarrow{\sim}\pi_{U,*}\Omega_{\mathcal C_U/U}
\]
(这说明 $\mathcal H$ 拟凝聚); 此外 $\pi_*\Omega_{\mathcal C/V}$ 是 $V$ 上秩 $g$ 的向量丛 (这说明 $\mathcal H$ 是秩 $g$ 的向量丛).
\end{example}
\begin{example}\label{src:5.1.21}
若 $G$ 是视作域 $k$ 上群概形的有限抽象群, 则 $BG$ 上的拟凝聚层对应于 $G$ 的表示 $V$. 若 $G$ 作用于仿射 $k$-概形 $\Spec A$, 则 $[\Spec A/G]$ 上的拟凝聚层由一个 $A$-模 $M$ 及群同态 $G\to\operatorname{End}_A(M)$ 给出. 这些描述由习题~\bookref{5.1.24}(b) 推出.
\end{example}
\begin{exercise}[等价性 I]\label{src:5.1.22}
设 $X$ 是概形.
\begin{enumerate}[(a)]
\item 证明对 $\mathcal O_{X_{\text{ét}}}$-模 $F$, 以下条件等价:
\begin{enumerate}[(i)]
\item $F$ 按定义~\bookref{5.1.17} 拟凝聚.
\item 存在拟凝聚 $\mathcal O_{X_{\mathrm{Zar}}}$-模 $F'$ 及同构 $F\cong\alpha^*F'$, 其中 $\alpha^*:\Mod(\mathcal O_{X_{\mathrm{Zar}}})\to\Mod(\mathcal O_{X_{\text{ét}}})$ 是将 $\mathcal O_{X_{\mathrm{Zar}}}$-模 $F'$ 送到 $\mathcal O_{X_{\text{ét}}}$-模 $\alpha^*F'$ 的函子, 后者定义为 $\alpha^*F'(U\xrightarrow{g}X)=\Gamma(U,g^*F')$.
\item $F$ 是赋环景 $(X_{\text{ét}},\mathcal O_{X_{\text{ét}}})$ 中的拟凝聚模对象, 这里使用一般赋环景 $(\mathcal S,\mathcal O)$ 中的定义~\cite[Tag 03DL]{SP}: $\mathcal O$-模 $F$ 称为拟凝聚的, 若对每个对象 $U\in\mathcal S$, 存在覆盖 $\{U_i\to U\}$, 使得它在限制景 $\mathcal S/U_i$ 上的限制 $F|_{U_i}$ 有自由呈示
\[
\mathcal O_{U_i}^{\oplus J}\longrightarrow\mathcal O_{U_i}^{\oplus I}\longrightarrow F|_{U_i}\longrightarrow0.
\]
\end{enumerate}
\textit{提示:} 证明 (iii) 蕴含其余条件时, 使用拟凝聚层的 fpqc 下降 (\bookref{3.1.4}) 构造拟凝聚层.
\item 证明 $\alpha^*$ 限制为范畴等价 $\QCoh(X_{\mathrm{Zar}})\to\QCoh(X_{\text{ét}})$.

\textit{提示:} 证明将 $\mathcal O_{X_{\text{ét}}}$-模限制到 Zariski 拓扑的运算 $F|_{X_{\mathrm{Zar}}}$ 定义了 $\alpha^*$ 的左伴随 $\alpha_*$. 利用伴随关系, 以及对拟凝聚 $\mathcal O_{X_{\mathrm{Zar}}}$-模 $F$ 的等同 $\alpha_*\alpha^*F\cong F$, 证明 $\alpha^*$ 在拟凝聚层上完全忠实.
\end{enumerate}
\end{exercise}
\begin{exercise}\label{src:5.1.23}
给出一个概形 $X$ 及非拟凝聚的 $\mathcal O_{X_{\text{ét}}}$-模 $F$, 使得 $F|_{X_{\mathrm{Zar}}}$ 拟凝聚.
\end{exercise}
这说明定义~\bookref{5.1.17} 中条件 (2) 的必要性.
\begin{exercise}[等价性 II]\label{src:5.1.24}
设 $\mathcal X$ 是 Deligne--Mumford 叠.
\begin{enumerate}[(a)]
\item 证明对 $\mathcal O_{\mathcal X_{\text{ét}}}$-模 $F$, 以下条件等价:
\begin{enumerate}[(i)]
\item $F$ 按定义~\bookref{5.1.17} 拟凝聚.
\item 存在来自概形的平展呈示 $U\to\mathcal X$, 使得 $F|_{U_{\text{ét}}}$ 拟凝聚.
\item $F$ 是赋环景 $(\mathcal X_{\text{ét}},\mathcal O_{\mathcal X_{\text{ét}}})$ 中的拟凝聚模层.
\end{enumerate}
\item 设 $U\to\mathcal X$ 是平展呈示, 置 $R=U\times_{\mathcal X}U$. 证明 $\mathcal X$ 上的拟凝聚层范畴等价于对 $(G,\alpha)$ 的范畴, 其中 $G$ 是 $U$ 上的拟凝聚层, $\alpha:p_1^*G\xrightarrow{\sim}p_2^*G$ 是 $R:=U\times_{\mathcal X}U$ 上满足上循环条件 $p_{23}^*\alpha\circ p_{12}^*\alpha=p_{13}^*\alpha$ 的同构.
\item 证明 $\mathcal X$ 上的拟凝聚层等价于如下数据: 对 $\mathcal X$ 上的每个概形 $U$, 给定拟凝聚层 $F_U$, 并对 $\mathcal X$ 上概形之间的每个映射 $f:U\to V$, 给定相容的同构 $f^*F_V\to F_U$.
\end{enumerate}
\end{exercise}
习题~\bookref{5.1.24}(c) 给出了理解模叠上拟凝聚层的一种方便且函子性的方式. 例如, $\mathcal M_g$ 上的 Hodge 丛可视为对每个光滑曲线族 $\pi:\mathcal C\to S$ 给定层 $\pi_*\Omega_{\mathcal C/S}$ 的数据.
\begin{exercise}[等价性 III]\label{src:5.1.25}
给定 Deligne--Mumford 叠 $\mathcal X$, 类似于 \S\bookref{3.2.2-subsection} 中概形上的定义, 定义大平展景 $(\Sch/\mathcal X)_{\text{ét}}$, 大 fppf 景 $(\Sch/\mathcal X)_{\mathrm{fppf}}$ 和光滑平展景 $\mathcal X_{\text{lis-ét}}$. 每个景都有结构层 $\mathcal O$, 称 $\mathcal O$-模拟凝聚, 若它作为赋环景中的模层拟凝聚. 证明存在等价
\[
\begin{aligned}
\QCoh(\mathcal X_{\text{ét}})&\xrightarrow{\sim}\QCoh(\mathcal X_{\text{lis-ét}})\xrightarrow{\sim}\QCoh((\Sch/\mathcal X)_{\text{ét}})\\
&\xrightarrow{\sim}\QCoh((\Sch/\mathcal X)_{\mathrm{fppf}}),
\end{aligned}
\]
其定义类似于习题~\bookref{5.1.22} 中概形 $X$ 上的等价 $\QCoh(X_{\mathrm{Zar}})\xrightarrow{\sim}\QCoh(X_{\text{ét}})$.
\end{exercise}

\subsection{推前与拉回}\label{src:5.1.4-subsection}
\begin{exercise}[推前--拉回伴随]\label{src:5.1.26}
设 $f:\mathcal X\to\mathcal Y$ 是 Deligne--Mumford 叠的态射.
\begin{enumerate}[(a)]
\item 证明 $f^*$ 保持拟凝聚性.
\item 若 $f$ 拟紧且拟分离, 证明 $f_*$ 保持拟凝聚性, 并推出
\[
\begin{tikzcd}
\QCoh(\mathcal X)\arrow[r,shift left,"f_*"]&\QCoh(\mathcal Y)\arrow[l,shift left,"f^*"]
\end{tikzcd}
\]
是伴随函子 ($f_*$ 为右伴随).
\end{enumerate}
\end{exercise}
\begin{exercise}[有益的练习]\label{src:5.1.27}
设 $G$ 是有限群, $k$ 是域.
\begin{enumerate}[(a)]
\item 考虑 $G$ 的表示 $V$, 并利用例~\bookref{5.1.21} 将它也视作 $BG$ 上的拟凝聚层. 对复合 $\Spec k\xrightarrow{p}BG\xrightarrow{\pi}\Spec k$, 证明 $\pi_*V=V^G$, 其中 $V^G$ 是 $G$-不变量子空间; 而 $p^*V=V$, 即忘掉 $G$-作用. 此外, 对 $k$-向量空间 $W$, 证明 $\pi^*W=W$, 赋予平凡 $G$-作用, 且 $p_*W=W\otimes p_*k$, 其中 $p_*k$ 是正则表示 $\Gamma(G,\mathcal O_G)$.
\item 给定 $G$ 在仿射 $k$-概形 $\Spec A$ 上的作用, 考虑图
\[
\begin{tikzcd}
\Spec A\arrow[r,"p"]&{[\Spec A/G]}\arrow[r,"\pi"]\arrow[d,"q"]&\Spec A^G\\
&BG&
\end{tikzcd}
\]
并回顾例~\bookref{5.1.21}, $[\Spec A/G]$ 上的拟凝聚层是带群同态 $G\to\operatorname{End}_A(M)$ 的 $A$-模 $M$. 显式描述拟凝聚层上的函子 $p_*,p^*,\pi_*,\pi^*,q_*,q^*$.
\end{enumerate}
\end{exercise}
\begin{exercise}[困难]\label{src:5.1.28}
设 $\mathcal X$ 是 Noether Deligne--Mumford 叠. 证明以下两条陈述:
\begin{enumerate}[(a)]
\item $\mathcal X$ 上的每个拟凝聚层都是其凝聚子层的有向余极限.
\item 若 $\mathcal U\subseteq\mathcal X$ 是开子叠, 则 $\mathcal U$ 上的每个凝聚层均可延伸为 $\mathcal X$ 上的凝聚层.
\end{enumerate}
\end{exercise}
这道习题将~\cite[习题 II.5.15]{Har77} 从概形推广到 Deligne--Mumford 叠; 另见~\cite[命题 15.4]{LMB00}, \cite[命题 7.1.11]{Ols16} 和~\cite[Tag 0GRE]{SP}.

\subsection{拟凝聚构造}\label{src:5.1.5-subsection}
Deligne--Mumford 叠上的\newterm{拟凝聚 $\mathcal O_{\mathcal X}$-代数}{quasi-coherent $\mathcal O_{\mathcal X}$-algebra} 是带有相容的 $\mathsf{Ab}(\mathcal X_{\text{ét}})$ 中环对象结构的拟凝聚 $\mathcal O_{\mathcal X}$-模. 定义\newterm{相对谱}{relative spectrum} $\operatorname{Spec}_{\mathcal X}\mathcal A$ 为如下叠: 它在概形 $S$ 上的对象由态射 $f:S\to\mathcal X$ 及 $\mathcal O_S$-代数态射 $f^*\mathcal A\to\mathcal O_S$ 组成.
\begin{exercise}[简单, 有益的练习]\label{src:5.1.29}
证明 $\operatorname{Spec}_{\mathcal X}\mathcal A$ 是 $\mathcal X$ 上的仿射代数叠.
\end{exercise}
\begin{example}[约化]\label{src:5.1.30}
设 $\mathcal X$ 是 Deligne--Mumford 叠, 以 $\mathcal O^{\mathrm{red}}_{\mathcal X}$ 记如下定义的 $\mathcal O_{\mathcal X}$-代数层: 对平展 $\mathcal X$-概形 $U$, 置 $\mathcal O^{\mathrm{red}}_{\mathcal X}(U\to\mathcal X)=\Gamma(U,\mathcal O_U)_{\mathrm{red}}$. 则 $\mathcal O^{\mathrm{red}}_{\mathcal X}$ 是拟凝聚 $\mathcal O_{\mathcal X}$-代数, 且 $\mathcal X_{\mathrm{red}}:=\operatorname{Spec}_{\mathcal X}\mathcal O^{\mathrm{red}}_{\mathcal X}$ 是约化代数叠, 带有闭浸入 $\mathcal X_{\mathrm{red}}\hookrightarrow\mathcal X$, 满足 $|\mathcal X_{\mathrm{red}}|=|\mathcal X|$. 称 $\mathcal X_{\mathrm{red}}$ 为 $\mathcal X$ 的\newterm{约化}{reduction}.
\end{example}
\begin{example}[正规化]\label{src:5.1.31}
设 $\mathcal X$ 是整 Deligne--Mumford 叠, 以 $\mathcal A$ 记 $\mathcal O_{\mathcal X}$-代数层, 其中 $\mathcal A(U\to\mathcal X)$ 定义为 $\Gamma(U,\mathcal O_U)$ 的正规化. 由于正规化与平展扩张交换 (命题~\bookref{A.7.4}), $\mathcal A$ 是拟凝聚 $\mathcal O_{\mathcal X}$-代数. $\mathcal X$ 的\newterm{正规化}{normalization} 定义为 $\widetilde{\mathcal X}:=\operatorname{Spec}_{\mathcal X}\mathcal A$.
\end{example}
\begin{exercise}[有益的练习]\label{src:5.1.32}
设 $f:\mathcal X\to\mathcal Y$ 是 Deligne--Mumford 叠的拟紧且拟分离态射.
\begin{enumerate}[(a)]
\item 证明存在分解 $f:\mathcal X\to\Spec f_*\mathcal O_{\mathcal X}\to\mathcal Y$.
\item 证明 $f$ 仿射当且仅当 $\mathcal X\to\Spec f_*\mathcal O_{\mathcal X}$ 是同构.
\item 证明 $f$ 拟仿射当且仅当 $\mathcal X\to\Spec f_*\mathcal O_{\mathcal X}$ 是开浸入.
\end{enumerate}
\end{exercise}
\begin{exercise}[细节]\label{src:5.1.33}
利用习题~\bookref{5.1.28} 证明 Noether Deligne--Mumford 叠上的每个拟凝聚代数层都是有限型子代数的有向余极限.
\end{exercise}

\subsection{上同调}\label{src:5.1.6-subsection}
我们建立 Deligne--Mumford 叠上 Abel 层的上同调理论. 虽然本讲义始终使用概形上拟凝聚层的上同调, 但令人意外的是, 我们很少需要代数空间和 Deligne--Mumford 叠上的上同调; 此处许多结果只是为了完整性而收入.

足够多内射对象的存在性可类似于概形的情形证明~\cite[命题 2.2]{Har77}.
\begin{lemma}\label{src:5.1.34}
若 $\mathcal X$ 是 Deligne--Mumford 叠, 则范畴 $\mathsf{Ab}(\mathcal X_{\text{ét}})$ 与 $\Mod(\mathcal O_{\mathcal X})$ 有足够多的内射对象. 若此外 $\mathcal X$ 拟分离, 则 $\QCoh(\mathcal X)$ 有足够多的内射对象.
\end{lemma}
\begin{proof}
回顾 Abel 范畴之间有正合左伴随 $L$ 的函子 $R:\mathcal A\to\mathcal B$ 保持内射对象, 因为对内射对象 $I\in\mathcal A$, $\operatorname{Hom}_{\mathcal B}(-,R(I))=\operatorname{Hom}_{\mathcal A}(L(-),I)$ 正合. 我们断言对 $\mathcal X_{\text{ét}}$ 上的每个环层 $\Lambda$, $\Lambda$-模范畴 $\Mod(\Lambda)$ 都有足够多的内射对象. 取 $\Lambda$ 为常值层 $\underline{\mathbb Z}$ (分别为结构层 $\mathcal O_{\mathcal X}$), 就说明了 $\mathsf{Ab}(\mathcal X_{\text{ét}})$ (分别为 $\Mod(\mathcal O_{\mathcal X})$) 有足够多的内射对象. 设 $F$ 是 $\Lambda$-模, $U\to\mathcal X$ 是平展呈示. 对每个几何点 $\Spec k\to U$, 考虑复合 $j_u:\Spec k\to U\to\mathcal X$. 由于 $\Spec k$ 上的平展层对应于普通的群 (习题~\bookref{5.1.5}), 逆像 $F_u:=j_u^{-1}F$ 对应于环 $\Lambda_u:=j_u^{-1}\Lambda$ 上的普通模. 选取到内射 $\Lambda_u$-模的嵌入 $F_u\hookrightarrow I_u$. 伴随给出映射 $F\to j_{u,*}I_u$, 且由于 $j_u^{-1}$ 正合 (习题~\bookref{5.1.10}), $j_{u,*}I_u\in\mathsf{Ab}(\mathcal X_{\text{ét}})$ 内射. 取乘积便得到到内射 $\Lambda$-模的单射 $F\hookrightarrow\prod_{u\in U}j_{u,*}I_u$.

对于最后一条陈述, 设 $F\in\QCoh(\mathcal X)$, $p:U=\coprod_i\Spec A_i\to\mathcal X$ 是平展呈示. 选取到内射拟凝聚 $\mathcal O_U$-模的单射 $p^*F\hookrightarrow I$. 复合 $F\hookrightarrow p_*p^*F\hookrightarrow p_*I$ 是单射, 且由于 $p^*$ 正合, $p_*I$ 内射.
\end{proof}
\begin{remark}\label{src:5.1.35}
上述论证可推广到有足够多点的任意赋环景 (见~\cite[定理 2.3.2]{Ols16}), 而结果甚至在任意赋环景上都成立~\cite[Tag 01DP]{SP}. Deligne--Mumford 叠 $\mathcal X$ 上的拟凝聚层范畴是 Grothendieck Abel 范畴~\cite[Tag 0781]{SP}, 每个这样的范畴都有足够多的内射对象~\cite{Gro57b}, \cite[Tag 079H]{SP}.
\end{remark}
\begin{definition}[上同调]\label{src:5.1.36}
设 $\mathcal X$ 是 Deligne--Mumford 叠, $F$ 是 $\mathcal X_{\text{ét}}$ 上的 Abel 群层. \newterm{上同调群}{cohomology group} $H^p(\mathcal X_{\text{ét}},F)$ 定义为整体截面函子 $\Gamma:\mathsf{Ab}(\mathcal X_{\text{ét}})\to\mathsf{Ab}$ 的第 $p$ 个右导出函子.

给定 Deligne--Mumford 叠的态射 $f:\mathcal X\to\mathcal Y$, \newterm{高阶直像}{higher direct image} $R^pf_*F$ 定义为 $f_*:\mathsf{Ab}(\mathcal X_{\text{ét}})\to\mathsf{Ab}(\mathcal Y_{\text{ét}})$ 的第 $p$ 个右导出函子.
\end{definition}
为了说明概形上拟凝聚层的上同调既可用通常的 Zariski 拓扑计算, 也可用小平展景计算 (命题~\bookref{5.1.41}), 我们将使用 Čech 上同调. 与概形情形一样, Čech 上同调提供了计算上同调群的极为有效的工具. 平展拓扑中的 Čech 上同调与 Zariski 拓扑中的情形~\cite[III.4]{Har77} 定义类似, 只需将交 $U_{i_0}\cap\cdots\cap U_{i_n}$ 换为纤维积 $U_{i_0}\times_{\mathcal X}\cdots\times_{\mathcal X}U_{i_n}$, 并考虑以任意顺序排列的所有 (可能重复的) 指标 $i_0,\ldots,i_n$.
\begin{definition}[Čech 上同调]\label{src:5.1.37}
设 $F$ 是 Deligne--Mumford 叠 $\mathcal X$ 上的 Abel 层. 若 $\mathcal U=\{U_i\to\mathcal X\}_{i\in I}$ 是平展覆盖, 则 $F$ 关于 $\mathcal U$ 的\newterm{Čech 复形}{Čech complex} 为 $\check C^\bullet(\mathcal U,F)$, 其中
\[
\check C^n(\mathcal U,F)=\prod_{(i_0,\ldots,i_n)\in I^{n+1}}F(U_{i_0}\times_{\mathcal X}\cdots\times_{\mathcal X}U_{i_n}),
\]
微分为
\[
\begin{aligned}
d^n:\check C^n(\mathcal U,F)&\longrightarrow\check C^{n+1}(\mathcal U,F),\\
(s_{i_0,\ldots,i_n})&\longmapsto\left(\sum_{k=0}^{n+1}(-1)^k\widehat p_k^*s_{i_0,\ldots,\widehat{i_k},\ldots,i_{n+1}}\right)_{(i_0,\ldots,i_{n+1})},
\end{aligned}
\]
其中
\[
\widehat p_k:U_{i_0}\times_{\mathcal X}\cdots\times_{\mathcal X}U_{i_{n+1}}\longrightarrow U_{i_0}\times_{\mathcal X}\cdots\times_{\mathcal X}\widehat{U_{i_k}}\times_{\mathcal X}\cdots\times_{\mathcal X}U_{i_{n+1}}
\]
是忘掉第 $k$ 个分量的映射 (指标从 $0$ 开始). $F$ 关于 $\mathcal U$ 的\newterm{Čech 上同调}{Čech cohomology} 为
\[
\check H^p(\mathcal U,F):=H^p(\check C^\bullet(\mathcal U,F)).
\]
\end{definition}
下面是 Čech 上同调中的标准结果, 其证明与拓扑空间情形类似. 它通常称为 Cartan 判据.
\begin{lemma}\label{src:5.1.38}
设 $\mathcal X$ 是 Deligne--Mumford 叠, $F$ 是 $\mathcal X_{\text{ét}}$ 上的 Abel 层. 设 $\mathrm{Ob}'$ 是 $\mathcal X$ 上平展概形集合的子集, $\mathrm{Cov}'$ 是 $\mathcal X_{\text{ét}}$ 中所有覆盖的子集, 满足对每个 $\{U_i\to U)\}\in\mathrm{Cov}'$, $U$, 各 $U_i$ 以及各纤维积 $U_{i_1}\times_U\cdots\times_U U_{i_n}$ 都属于 $\mathrm{Ob}'$. 假设平展 $\mathcal X$-概形 $U$ 的每个覆盖在 $\mathrm{Cov}'$ 中都有加细. 若对每个覆盖 $\mathcal U\in\mathrm{Cov}'$ 及 $p>0$ 都有 $\check H^p(\mathcal U,F)=0$, 则对每个 $U\in\mathrm{Ob}'$ 有 $H^p(U,F)=0$.
\end{lemma}
\begin{proof}
这在任意景中成立; 见~\cite[II.5.9.2]{God58}, \cite[命题 III.2.12]{Mil80}, \cite[Tag 03F9]{SP} 或~\cite[命题 2.3.15]{Ols16}.
\end{proof}
与普通拓扑空间的情形~\cite[习题 III.4.11]{Har77} 一样, 平展层的上同调可作为关于一个上同调消失的覆盖的 Čech 上同调来计算.
\begin{lemma}\label{src:5.1.39}
设 $F$ 是 $\mathcal X_{\text{ét}}$ 上的 Abel 层, $(U_i\to\mathcal X)_{i\in I}$ 是平展覆盖. 若对所有 $p>0$, $n\geq0$ 及 $i_0,\ldots,i_n\in I$ 都有 $H^p(U_{i_0}\times_U\cdots\times_U U_{i_n},F)=0$, 则 $\check H^p(\mathcal U,F)=H^p(\mathcal X_{\text{ét}},F)$.
\end{lemma}
\begin{proof}
这也在任意景中成立, 并可由 Čech 到导出函子的谱序列推出; 见~\cite[Tag 03F7]{SP}.
\end{proof}
\begin{theorem}\label{src:5.1.40}
对仿射概形 $X$ 上的拟凝聚层 $F$, 对所有 $p>0$ 都有 $H^p(X_{\text{ét}},F)=0$.
\end{theorem}
\begin{proof}
设 $X=\Spec A$, $F=\widetilde M$ 是拟凝聚 $\mathcal O_{X_{\mathrm{Zar}}}$-模, 它也可视作拟凝聚 $\mathcal O_{X_{\text{ét}}}$-模 (习题~\bookref{5.1.22}). 形如 $\mathcal U=\{\Spec B\to\Spec A\}$ 的平展覆盖足以加细任意其他覆盖. 对覆盖 $\mathcal U$, 忠实平坦下降 (习题~\bookref{3.1.3}) 给出长正合列
\[
0\longrightarrow M\longrightarrow M\otimes_A B\longrightarrow M\otimes_A B\otimes_A B\longrightarrow M\otimes_A B\otimes_A B\otimes_A B\longrightarrow\cdots,
\]
它可与 Čech 复形 $\check C^\bullet(\mathcal U,F)$ 等同. 这说明对 $p>0$ 有 $\check H^p(\mathcal U,F)=0$, 因而引理~\bookref{5.1.38} 蕴含 $H^p(X_{\text{ét}},F)=0$.
\end{proof}
\begin{proposition}\label{src:5.1.41}
若 $X$ 是具有仿射对角线的概形, $F$ 是拟凝聚层, 则对所有 $p$ 都有 $H^p(X,F)=H^p(X_{\text{ét}},F)$.
\end{proposition}
\begin{proof}
设 $\mathcal U:=\{U_i\hookrightarrow X\}$ 是由仿射概形组成的 $X$ 的 Zariski 覆盖. 由于 $X$ 具有仿射对角线, 每个纤维积 $U_{i_0}\times_X\cdots\times_X U_{i_n}$ 也仿射. 由定理~\bookref{5.1.40}, $H^p(U_{i_0}\times_U\cdots\times_U U_{i_n},F)=0$, 故可应用引理~\bookref{5.1.39}, 推出 $H^p(X,F)=\check H^p(\mathcal U,F)=H^p(X_{\text{ét}},F)$. 同样的结果更一般地在不假设仿射对角线时, 以及在其他景中, 例如大 fppf 景 $(\Sch/X)_{\mathrm{fppf}}$ 中成立; 见~\cite[命题 3.7]{Mil80} 和~\cite[Tag 03DW]{SP}.
\end{proof}
现在比较在 $\mathsf{Ab}(\mathcal X_{\text{ét}})$, $\Mod(\mathcal O_{\mathcal X_{\text{ét}}})$ 与 $\QCoh(\mathcal X)$ 中计算的上同调.
\begin{proposition}\label{src:5.1.42}
设 $\mathcal X$ 是 Deligne--Mumford 叠.
\begin{enumerate}[(1)]
\item 若 $F$ 是 $\mathcal O_{\mathcal X}$-模, 则 $F$ 作为 Abel 层的上同调 $H^p(\mathcal X_{\text{ét}},F)$ 与 $\Gamma:\Mod(\mathcal O_{\mathcal X})\to\mathsf{Ab}$ 的第 $p$ 个右导出函子一致.
\item 若 $\mathcal X$ 拟紧且具有仿射对角线, $F$ 是 $\mathcal X$ 上的拟凝聚层, 则 $F$ 作为 Abel 层的上同调 $H^p(\mathcal X_{\text{ét}},F)$ 与 $\Gamma:\QCoh(\mathcal X)\to\mathsf{Ab}$ 的第 $p$ 个右导出函子一致.
\end{enumerate}
对 Deligne--Mumford 叠的态射 $f:\mathcal X\to\mathcal Y$ 的高阶直像 $R^pf_*F$ 也成立.
\end{proposition}
\begin{proof}
对于 (1), 需要说明 $\Mod(\mathcal O_{\mathcal X})$ 中的内射对象在 $\mathsf{Ab}(\mathcal X_{\text{ét}})$ 中无上同调. 这里使用 Čech 上同调中的标准技巧. 约定如下记号: 若 $\mathcal U=\{U_i\to\mathcal X\}_{i\in I}$ 是平展覆盖, $j=(j_0,\ldots,j_n)\in I^{n+1}$, 则 $U_j:=U_{j_0}\times_{\mathcal X}\cdots\times_{\mathcal X}U_{j_n}$, 其结构态射为 $q_j:U_j\to\mathcal X$. 在 $\mathcal X$ 上有预层的链复形 $\underline{\mathbb Z}_{\mathcal U,\bullet}$, 定义为
\[
\underline{\mathbb Z}_{\mathcal U,n}:=\bigoplus_{j\in I^{n+1}}q_{j,!}\underline{\mathbb Z},
\]
其中 $\underline{\mathbb Z}$ 表示常值预层, 且 $q_{j,!}\underline{\mathbb Z}(V\to\mathcal X):=\bigoplus_{\operatorname{Mor}_{\mathcal X}(V,U_i)}\underline{\mathbb Z}$, 微分为自然映射的交错和. 此预层复形在正次数正合, 并且对每个预层 $F$ 都有
\[
\begin{aligned}
\check C(\mathcal U,F)&=\operatorname{Mor}_{\mathsf{PAb}(\mathcal X_{\text{ét}})}(\underline{\mathbb Z}_{\mathcal U,\bullet},F)\\
&=\operatorname{Mor}_{\mathsf{PMod}(\mathcal O_{\mathcal X})}(\underline{\mathbb Z}_{\mathcal U,\bullet}\otimes_{\mathbb Z}\mathcal O_{\mathcal X},F),
\end{aligned}
\]
其中态射在预层范畴 $\mathsf{PAb}(\mathcal X_{\text{ét}})$ 与 $\mathsf{PMod}(\mathcal O_{\mathcal X})$ 中计算. 若 $F\in\Mod(\mathcal O_{\mathcal X})$ 内射, 则它作为 $\mathcal O_{\mathcal X}$-模预层也内射. 因而 $\check C(\mathcal U,F)$ 在正次数正合, 故对 $p>0$ 有 $\check H^p(\mathcal U,F)=0$. 于是引理~\bookref{5.1.38} 蕴含 $H^p(\mathcal X_{\text{ét}},F)=0$. 此结果更一般地在任意赋环景中成立; 见~\cite[推论 2.3.16]{Ols16} 或~\cite[Tag 03FD]{SP}.

对于 (2), 设 $F\in\QCoh(\mathcal X)$ 是内射对象. 设 $q:U=\Spec A\to\mathcal X$ 是平展呈示, 选取到内射对象 $G\in\QCoh(U)$ 的单射 $q^*F\hookrightarrow G$. 则推前 $q_*G$ 内射 (因为右伴随 $q^*$ 正合), 且有内射对象之间的嵌入 $F\hookrightarrow q_*q^*F\hookrightarrow q_*G$, 它必分裂. 因此只需说明 $q_*G$ 在 $\mathsf{Ab}(\mathcal X_{\text{ét}})$ 中无上同调. 由于 $\mathcal X$ 具有仿射对角线, $p:U\to\mathcal X$ 是仿射态射. 由平坦基变换 (\bookref{A.2.12}) 及平展下降, $q_*$ 在拟凝聚层范畴上正合. 由定理~\bookref{5.1.40}, $H^p(\mathcal X_{\text{ét}},q_*G)=H^p(U_{\text{ét}},G)=0$.
\end{proof}
\begin{exercise}[平坦基变换]\label{src:5.1.43}
考虑 Deligne--Mumford 叠的 Cartesian 图
\[
\begin{tikzcd}
\mathcal X'\arrow[r,"{g'}"]\arrow[d,"{f'}"']&\mathcal X\arrow[d,"f"]\\
\mathcal Y'\arrow[r,"g"']\arrow[ur,phantom,"\square" description]&\mathcal Y,
\end{tikzcd}
\]
设 $F$ 是 $\mathcal X$ 上的拟凝聚层. 若 $g:\mathcal Y'\to\mathcal Y$ 平坦, $f:\mathcal X\to\mathcal Y$ 拟紧且拟分离, 则自然伴随映射
\[
g^*R^pf_*F\longrightarrow R^pf'_*g'^*F
\]
对所有 $p\geq0$ 都是同构.
\end{exercise}
\begin{exercise}\label{src:5.1.44}
若 $\mathcal X$ 是 Deligne--Mumford 叠, $F_i$ 是 $\mathcal X_{\text{ét}}$ 中 Abel 层的有向系, 证明 $\mathrm{colim}_i H^p(\mathcal X_{\text{ét}},F_i)\to H^p(\mathcal X_{\text{ét}},\mathrm{colim}_i F_i)$ 是同构.
\end{exercise}
\begin{example}[群上同调]\label{src:5.1.45}
设 $G$ 是视作域 $k$ 上群概形的有限抽象群, $V$ 是 $G$ 的表示, 也可视作 $BG$ 上的拟凝聚层. \newterm{群上同调}{group cohomology} $H^p(G,V)$ 定义为函子
\[
\operatorname{Rep}(G)\longrightarrow\operatorname{Vect}_k,\qquad V\longmapsto V^G
\]
的第 $p$ 个右导出函子. 由于 $\operatorname{Rep}(G)\cong\QCoh(BG)$, $\Gamma(BG,V)=V^G$, 且 $H^p((BG)_{\text{ét}},V)$ 可在 $\QCoh(BG)$ 中计算 (命题~\bookref{5.1.42}(2)), 有等同
\[
H^p(G,V)\cong H^p((BG)_{\text{ét}},V).
\]
$V$ 关于平展覆盖 $\mathcal U=\{\Spec k\to BG\}$ 的 Čech 复形 $\check C^\bullet(\mathcal U,V)$ 的各项为
\[
\check C^n(\mathcal U,V)=\Gamma(G,\mathcal O_G)^{\otimes n}\otimes V,
\]
微分为
\begin{equation}\label{src:5.1.46}
\begin{aligned}
d^n:\Gamma(G,\mathcal O_G)^{\otimes n}\otimes V&\longrightarrow\Gamma(G,\mathcal O_G)^{\otimes(n+1)}\otimes V,\\
f\otimes v&\longmapsto\sum_{k=0}^n(-1)^k\mu_k^*(f)\otimes v+(-1)^{n+1}f\otimes\sigma(v),
\end{aligned}
\end{equation}
其中 $\sigma:V\to\Gamma(G,\mathcal O_G)\otimes V$ 是余作用, $\mu_k:G^{n+1}\to G^n$ 定义为当 $k=0$ 时将 $(g_1,\ldots,g_{n+1})$ 送到 $(g_2,\ldots,g_{n+1})$, 而当 $k=1,\ldots,n$ 时送到 $(g_1,\ldots,g_{k-1},g_kg_{k+1},g_{k+2},\ldots,g_{n+1})$; 见习题~\bookref{4.4.6}. 在低次数, 有 $d^0(v)=v-\sigma(v)$ 及
\[
d^1(f,v)=1\otimes f\otimes v-\mu^*(f_1)\otimes v+f_1\otimes\sigma(v),
\]
其中 $\mu=\mu_1$ 是群乘法 $G\times G\to G$. 这称为\newterm{标准消解}{standard resolution} 或\newterm{bar 消解}{bar resolution}, 提供了计算群上同调的有效方法. 由于 $G$ 有限, 此复形也可利用与集合映射之间的等同 $\Gamma(G^n,\mathcal O_{G^n})\otimes V\cong\operatorname{Map}(G^n,V)$ 来解释. 另见~\cite[\S6]{Wei94} 和~\cite[\S1]{Bro94}.
\end{example}
以下例子说明 Deligne--Mumford 叠的凝聚层上同调在任意高次数都可能非零. 在特征 $0$ 或温和 Deligne--Mumford 叠的情形, 这不会发生 (见推论~\bookref{5.4.23}).
\begin{exercise}\label{src:5.1.47}
设 $k$ 是特征 $p>0$ 的域, $\mathbb Z/p\mathbb Z$ 是 $\Spec k$ 上对应于抽象循环群的群概形. 证明对每个 $i\geq0$ 都有 $H^i(B(\mathbb Z/p\mathbb Z),\mathcal O_{B(\mathbb Z/p\mathbb Z)})=k$.
\end{exercise}
\begin{exercise}[困难]\label{src:5.1.48}
将 Grothendieck 消失定理~\cite[定理 III.2.7]{Har77} 推广到代数空间: 若 $X$ 是 Noether 代数空间, $F$ 是拟凝聚层, 则对 $i>\dim X$ 有 $H^i(X,F)=0$. 另见~\cite[Tag 0A4Q]{SP}.
\end{exercise}

\subsection{Picard 群}\label{src:5.1.7-subsection}
以 $\operatorname{Pic}(\mathcal X)$ 记 Deligne--Mumford 叠 $\mathcal X$ 上线丛同构类的集合. 它在张量积下构成 Abel 群.
\begin{exercise}\label{src:5.1.49}
设 $k$ 是域, $n$ 是与 $\operatorname{char}(k)$ 互素的正整数. 这保证了 $B\mu_n$ 是 Deligne--Mumford 叠. 利用 Čech 上同调证明 $\operatorname{Pic}(B\mu_n)=\mathbb Z/n\mathbb Z$. (另一种方法及 $n$ 与 $\operatorname{char}(k)$ 不互素时的推广见例~\bookref{7.1.18}.)
\end{exercise}
\begin{exercise}[困难]\label{src:5.1.50}
证明在 $\operatorname{char}(k)\neq2,3$ 的域上, $\operatorname{Pic}(\mathcal M_{1,1})=\mathbb Z/12$, 且由 Hodge 丛生成 (见例~\bookref{5.1.4}).
\end{exercise}
这是 Mumford 影响深远的论文~\cite{Mum65} 的主要结论, 其中结果的陈述没有使用 Deligne--Mumford 叠的语言 (当时尚未引入). 有了更多工具, 即拟凝聚层之后, 这道习题会更容易; 见习题~\bookref{7.1.21}.
